\documentclass{article}
\usepackage{graphicx} % Required for inserting images

\title{FDTD: diffraction over objects}
\author{Alichan Borokov, Sam Hemati, Joran Van Haelst}
\date{December 2025}

\begin{document}
\maketitle

\section{Introduction}


\section{Infinitely Thin Structure}

\subsection{Analytical diffraction around wedge}

\subsection{Results}


\section{Thick Object with Finite Impedance}

\subsection{Results}


\section{Tilted Boundaries}

\subsection{Finite Impedance}

\subsection{Refinement}

In order to better model the non-cartesian surface,
we included a grid refinement around the tilted boundaries.

This was achieved by making \(dx\) and \(dy\) functions of respectively \(x\) and \(y\).
As a first attempt both were made a step-function,
in which around the boundaries they are both half of what they are elsewhere.
This however results in reflections at the discontinuity in \(dx\) or \(dy\), 
as can be seen in ??. % TODO: insert figure of reflections by step function in dx/dy
Next, to remove these reflections, 
a buffer zone was added in both \(x\)- and \(y\)-direction in which \(dx\) and \(dy\) are linearly interpolated (see fig ??). % TODO: add figure of dx or dy

Both these were in practice implemented by taking the original discretisation coordinates,
and transforming these, either partswise linearly (for the step-function in \(dx\) or \(dy\)),
or with a parabolic part in the bufferzone (for the linear \(dx\) and \(dy\)).

There were \(dx\) or \(dy\) is not constant however, 
this results in an error in the update scheme.
As, for example, the discretised \(o_x\) may no longer be halfway between two \(p\)'s.
This means that in the update equations of \(o_x\),
it is erronous to use \(\frac{dp}{dx}\) as calculated by differentiating two neighbouring \(p\)'s.
So an interpolation was made between three values of \(\frac{dp}{dx}\) at the halfway point between two \(p\)'s: 
two closest neightbours to the left, 
the two direct neighbours, and those to the right.
Since there are three values, 
a parabolic interpolation was optimal.
At the edges, 
where there are only two values to interpolate between,
a linear interpolation was used instead.
To update \(o_x\), 
the interpolated \(\frac{dp}{dx}\) at the location of \(o_x\) (after the coordinate transformation) was used.

Some results from the refinement and comparison with without: INSERT HERE. % TODO

\subsection{Non Cartesian Grid}

Another route that was explored is to tilt the grid.
In this way the surface of the tilted boundary becomes flat on one side,
thus removing the problems that arise from a staircase-like boundary.
The left-side was chosen to be flat, as the most important reflections occur there.

Instead of working with \(x\)- and \(y\)-coordinates, 
a switch is made to \(u\)- and \(v\)-coordinates. 
These are related via equations \ref{eq:x-of-u-and-v} and \ref{eq:y-of-u-and-v}, 
in which \(\theta\) is the angle between the \(v\)-axis and the \(x\)-axis.

\begin{equation}\label{eq:x-of-u-and-v}
    x = u + v \cos \theta
\end{equation}
\begin{equation}\label{eq:y-of-u-and-v}
    y = v \sin \theta
\end{equation}

In figure \ref{fig:x-y-w-and-v} the basis-vectors are given. % TODO: figure of ex, ey, ev and ew with annotation of theta
Note that \(\vec{e_x} = \vec{e_u}\), but \(x \ne u\). 
In this figure another vector, \(\vec{e_w}\), is introduced.
This vector is perpendicular to \(\vec{e_v}\).
With this figure the relations given in equations \ref{eq:ew-of-ex-and-ey} and \ref{eq:ev-of-ex-and-ey} are easily derived.

\begin{equation}\label{eq:ew-of-ex-and-ey}
    \vec{e_w} = \sin \theta \vec{e_x} - \cos \theta \vec{e_y}
\end{equation}
\begin{equation}\label{eq:ev-of-ex-and-ey}
    \vec{e_v} = \cos \theta \vec{e_x} + \sin \theta \vec{e_y}
\end{equation}

There is still an important choice remaining,
namely, which directions of \(\vec{o}\) to store.
A first idea might be to choose \(o_u = \vec{e_u} \cdot \vec{o}\) and \(o_v = \vec{e_v} \cdot \vec{o}\),
as analogous to the cartesian grid.
Remember however that the goal is to implement a tilted PEC-boundary,
which has as boundary condition that \(o_n = 0\), 
with \(n\) the normal of the boundary.
For this boundary condition to be implemented simply, 
\(o_w\) and \(o_v\) are chosen, 
in which \(\vec{e_w}\) is the normal of the tilted boundary (on the left side).
An added benefit is that now both stored directions of \(\vec{o}\) are orthogonal to eachother,
so no data overlaps. Figure \ref{fig:discretised-p-ow-and-ov} shows the location of a few \(p\), \(o_w\) and \(o_v\)'s. % TODO: figure showcasing the discretisation of p, ow and ov

Now the update scheme can be adjusted for the tilted grid.
To update \(p\), 
the divergence of \(\vec{o}\) is needed at the location of \(p\),
with \(\nabla \cdot \vec{o} = \frac{do_w}{dw} + \frac{do_v}{dv}\).
The following quantities are known:
\[ \frac{do_v}{dv}|_{v_0 + \Delta v/2} \cdot \Delta v \approx o_v|_{v_0 + \Delta v} - o_v|_{v_0} , \]
\[ \frac{do_w}{du}|_{u_0 + \Delta u/2} \cdot \Delta u \approx o_w|_{u_0 + \Delta u} - o_w|_{u_0} . \]
\(\frac{do_w}{dw}\) requires some more attention.
From equations \ref{eq:ev-of-ex-and-ey} and \ref{eq:ew-of-ex-and-ey}, 
the following can be derived:
\[ \vec{e_w} = \frac{1}{\sin \theta} \left( \vec{e_u} - \cos \theta \vec{e_v} \right) . \]
From which 
\[ \frac{do_w}{dw} = \vec{e_w} \cdot \nabla o_w = \frac{1}{\sin \theta} \left( \frac{do_w}{du} - \cos \theta \frac{do_w}{dv} \right) \] 
follows.
So the missing piece is \(\frac{do_w}{dv}\), 
which can only be interpolated at \(p\)'s location by using
\[ \frac{do_w}{dv}|_{v_0} \cdot 2 \Delta v \approx o_w|_{v_0 + \Delta v} - o_w|_{v_0 - \Delta v} , \]
combined with
\[ o_w|_{u_0 + \Delta u/2} \approx \frac{o_w|_{u_0} + o_w|_{u_0 + \Delta u}}{2} . \]
For the update equations of \(\vec{o}\) an analogous approach is followed.

\begin{equation}\label{eq:update-p-perp-PML}
    \frac{dp_\perp}{dt} + c^2 \frac{do_\alpha}{d\alpha} + \kappa_\perp p_\perp = 0
\end{equation}
\begin{equation}\label{eq:update-p-para-PML}
    \frac{dp_\parallel}{dt} + c^2 \frac{do_\gamma}{d\gamma} + \kappa_\parallel p_\parallel = 0
\end{equation}

The perfectly matched layers in the tilted grid become a bit more complex.
This especially for the upper PML.
The update equations for \(p = p_\perp + p_\parallel\) are given by equations \ref{eq:update-p-perp-PML} and \ref{eq:update-p-para-PML}.
Looking at the upper PML, 
so that the perpendicular direction is \(\vec{e_y}\) and the parallel one is \(\vec{e_x}\),
these equations become:
\[ \frac{dp_\perp}{dt} + c^2 \frac{do_y}{dy} + \kappa_\perp p_\perp = 0 , \]
\[ \frac{dp_\parallel}{dt} + c^2 \frac{do_x}{dx} + \kappa_\parallel p_\parallel = 0 . \]
By defining \(\vec{p} = p_\perp \vec{e_y} + p_\parallel \vec{e_x}\) and \(\vec{\kappa} = \kappa_\perp \vec{e_y} + \kappa_\parallel \vec{e_x}\),
these can be rewritten in one equation as follows
\[ \frac{d}{dt} \vec{e_i} \cdot \vec{p} + c^2 \vec{e_i} \cdot \nabla \vec{o} \cdot \vec{e_i} + \left| \vec{e_i} \cdot \vec{\kappa} \right| \vec{p} \cdot \vec{e_i} = 0\]
The absolute value of \(\vec{e_i} \cdot \vec{\kappa}\) is important: 
if we chose \(-\vec{e_y}\) instead of \(\vec{e_y}\) as the perpendicular direction, ... % TODO

CN

BOUNDARY


\section{Conclusion}

\end{document}
